%% GENERATED by tools/gen_latex.py -- edit content/ch07.py instead.
%% Build:  latexmk -pdf ch07.tex     (or: pdflatex ch07.tex, twice)
\documentclass[aspectratio=169,11pt,t]{beamer}
\usepackage{itbpro}

\coursetitle{Menyelam ke Keras}
\coursesubtitle{Tiga cara membangun model, tiga tingkat kendali atas pelatihannya -- dan satu asas yang menyatukan semuanya: mudah dimulai, tanpa langit-langit.}
\coursekicker{Chapter 7}
\coursesource{Chollet \& Watson, Deep Learning with Python 3e -- bab 7 (hlm. 190-230)}
\courseduration{3 jam (2 sesi)}
\coursepresenter{Rahman Indra Kesuma, S.Kom., M.Cs.}{Asisten Pengajar}
\coursebrand{AI for Professional \textperiodcentered\ ITB \textperiodcentered\ Ch 7}

\title{Menyelam ke Keras}
\author{Rahman Indra Kesuma, S.Kom., M.Cs.}
\date{}

\begin{document}
\covertitle

\framekicker{Learning outcomes}
\begin{frame}[shrink=25]{Tujuan Sesi}
\leadpar{Setelah sesi ini, peserta dapat:}
\begin{isteps}
\item Menjelaskan asas \textbf{progressive disclosure of complexity} dan menempatkan diri sendiri pada spektrum alur kerja Keras.
\item Membangun model bermasukan-jamak dan berkeluaran-jamak dengan \textbf{Functional API}, lalu melatihnya dengan senarai maupun kamus.
\item Memanfaatkan \textbf{akses ke keterhubungan lapis}: menggambar topologi dengan \inlinecode{plot\_model()} dan \textbf{mengekstraksi fitur} dari simpul antara.
\item Menulis \textbf{Model subclass} -- dan menyebut apa yang hilang saat memilihnya.
\item Menulis \textbf{metrik kustom} (\inlinecode{update\_state}, \inlinecode{result}, \inlinecode{reset\_state}) dan \textbf{callback kustom}, serta memakai \inlinecode{EarlyStopping} + \inlinecode{ModelCheckpoint}.
\item Menulis \textbf{\inlinecode{train\_step()} sendiri} di TensorFlow, PyTorch, dan JAX -- dan menyisipkannya ke \inlinecode{fit()} supaya callback dan optimasi bawaan tetap terpakai.
\end{isteps}

\vskip 6pt
\vskip 4pt
\reslink{Buku}{Bab 7 --- teks penuh}{https://deeplearningwithpython.io/chapters/chapter07\_deep-dive-keras}
\reslink{Notebook}{01\_tiga\_cara\_membangun\_model.ipynb}{../../notebooks/ch07/01\_tiga\_cara\_membangun\_model.ipynb}
\reslink{Notebook}{02\_metrik\_dan\_callback\_kustom.ipynb}{../../notebooks/ch07/02\_metrik\_dan\_callback\_kustom.ipynb}
\reslink{Notebook}{03\_train\_step\_kustom\_per\_backend.ipynb}{../../notebooks/ch07/03\_train\_step\_kustom\_per\_backend.ipynb}
\reslink{Repo}{Notebook resmi penulis}{https://github.com/fchollet/deep-learning-with-python-notebooks/blob/master/chapter07\_deep-dive-keras.ipynb}
\vskip 2pt
\end{frame}

\framekicker{Bagian 7.1}
\begin{frame}[shrink=25]{Progressive disclosure of complexity}
\begin{center}
\adjustbox{max width=\linewidth,max totalheight=0.52\textheight}{%

\begin{tikzpicture}[font=\sffamily\tiny,
  bx/.style={draw=rule, fill=papertint, rounded corners=4pt, minimum width=2.5cm,
             minimum height=1.25cm, align=left, text=ink2},
  ax/.style={draw=signal!60, fill=signal!9, rounded corners=4pt, minimum width=2.5cm,
             minimum height=1.25cm, align=left, text=ink2}]
  \shade[left color=itbbluelt!45, right color=violet!45]
    (0,1.15) rectangle (10.4,1.28);
  \node[bx] (a) at (1.25,0.4)
    {\textbf{Sequential API}\\+ lapis bawaan\\[2pt]\textcolor{ink3}{pemula, model sederhana}};
  \node[ax] (b) at (3.9,0.4)
    {\textbf{Functional API}\\+ lapis bawaan\\[2pt]\textcolor{ink3}{kasus pemakaian baku}};
  \node[ax] (c) at (6.55,0.4)
    {\textbf{Functional API}\\+ lapis, metrik, rugi kustom\\[2pt]\textcolor{ink3}{kasus khusus}};
  \node[draw=violet!70, fill=violet!12, rounded corners=4pt, minimum width=2.5cm,
        minimum height=1.25cm, align=left, text=ink2] (d) at (9.2,0.4)
    {\textbf{Subclassing}\\tulis semuanya sendiri\\[2pt]\textcolor{ink3}{peneliti}};
  \node[text=signal, anchor=west] at (0,-0.5) {mudah dipakai};
  \node[text=violet, anchor=east] at (10.4,-0.5) {luwes sepenuhnya};
  \node[draw=signal!35, fill=signal!6, rounded corners=4pt, minimum width=10.2cm,
        minimum height=0.7cm, align=left, text=ink2] at (5.2,-1.15)
    {~Semuanya berdiri di atas API bersama: \texttt{Layer} dan \texttt{Model}. Komponen dari satu\\
     ~alur kerja bisa dipakai di alur kerja mana pun --- tak perlu ganti kerangka saat naik tingkat.};
\end{tikzpicture}

}
\end{center}
\vskip 3pt{\color{ink3}\fontsize{7.4}{9.4}\selectfont Gambar 7.1 -- bukan empat kerangka berbeda, melainkan satu spektrum di atas API yang sama.\par}
\begin{iquote}
\quotetext{Anda bisa memakai Keras seperti memakai scikit-learn -- tinggal memanggil \inlinecode{fit()} dan membiarkan kerangkanya bekerja -- atau memakainya seperti NumPy, dengan kendali penuh atas setiap detail kecil.}\quotecite{Chollet \& Watson, bab 7.1}
\end{iquote}
\begin{iband}[signal]
Analogi yang dipakai buku: \textbf{Keras adalah Python-nya deep learning}. Python multiparadigma -- berorientasi objek, fungsional, prosedural, semuanya rukun. Karena itu \hilite{Anda tidak perlu pindah kerangka} saat berpindah dari mahasiswa ke peneliti, atau dari data scientist ke deep learning engineer.
\end{iband}
\end{frame}
\note{Kalimat yang perlu diulang: apa yang dipelajari peserta hari ini tetap berlaku setelah mereka jadi ahli. Ini menghilangkan kecemasan 'nanti harus belajar ulang'.}

\coursesection{01}{Tiga cara membangun model}{Sequential, Functional, subclassing -- dan kapan masing-masing.}

\framekicker{Bagian 7.2.1-7.2.2}
\begin{frame}[shrink=25]{Sequential itu pada dasarnya senarai Python}
\begin{minipage}[t]{0.485\textwidth}
{\fontsize{9.4}{12.6}\selectfont\color{fg2}\textbf{Sequential} -- API yang paling mudah didekati. Batasnya juga jelas: ia hanya bisa menyatakan model dengan \textbf{satu masukan dan satu keluaran}, satu lapis sesudah lapis lain.\par}\vskip 4pt
\begin{iband}[amber]
Padahal di praktik, model dengan \textbf{masukan jamak} (citra dan metadatanya), \textbf{keluaran jamak} (beberapa hal yang mau diramalkan), atau \textbf{topologi tak-linear} itu \hilite{biasa sekali}
\end{iband}

\end{minipage}%
\hfill
\begin{minipage}[t]{0.485\textwidth}
\icode{python}{listing 7.8 --- versi Functional}{listings/ch07-01.py}

\end{minipage}%
\icode{python}{tensor simbolik}{listings/ch07-02.py}
\begin{iband}[signal]
\inlinecode{inputs} itu \textbf{tensor simbolik}: ia \hilite{tidak memuat data apa pun}, tetapi menyandikan spesifikasi tensor yang kelak akan dilihat model. Semua lapis Keras bisa dipanggil baik pada tensor data sungguhan maupun pada tensor simbolik -- yang kedua mengembalikan tensor simbolik baru dengan shape dan dtype yang sudah diperbarui.
\end{iband}
\end{frame}

\framekicker{Bagian 7.2.2 $\cdot$ listing 7.9}
\begin{frame}[shrink=25]{Di sinilah Functional API bersinar}
{\fontsize{9.4}{12.6}\selectfont\color{fg2}Kasus dari buku: sistem yang \textbf{memeringkat tiket dukungan pelanggan menurut prioritas} dan mengarahkannya ke departemen yang tepat.\par}\vskip 4pt
\begin{center}
\adjustbox{max width=\linewidth,max totalheight=0.52\textheight}{%

\begin{tikzpicture}[font=\sffamily\tiny,
  ix/.style={draw=rule, fill=papertint, rounded corners=3pt, minimum width=1.9cm,
             minimum height=0.5cm, text=ink2, font=\ttfamily\tiny},
  ax/.style={draw=signal!60, fill=signal!10, rounded corners=3pt, minimum width=1.9cm,
             minimum height=0.68cm, text=ink, align=center},
  ox/.style={draw=lime!65, fill=limebr!14, rounded corners=3pt, minimum width=1.7cm,
             minimum height=0.6cm, text=ink, align=center},
  ar/.style={-{Stealth[length=4pt]}, signal, line width=0.7pt}]
  \node[ix] (t)  at (0,1.1)  {title};
  \node[ix] (b)  at (0,0.35) {text\_body};
  \node[ix] (g)  at (0,-0.4) {tags};
  \node[ax] (cc) at (2.7,0.35) {Concatenate\\\ttfamily (20100,)};
  \node[draw=violet!70, fill=violet!12, rounded corners=3pt, minimum width=2.1cm,
        minimum height=0.68cm, text=ink, align=center] (d) at (5.5,0.35)
        {Dense 64 relu\\\ttfamily dense\_features};
  \node[ox] (p)  at (8.2,0.95) {priority\\\ttfamily 1 $\cdot$ sigmoid};
  \node[ox] (dp) at (8.2,-0.3) {department\\\ttfamily 4 $\cdot$ softmax};
  \draw[ar] (t) -- (cc); \draw[ar] (b) -- (cc); \draw[ar] (g) -- (cc);
  \draw[ar] (cc) -- (d); \draw[ar] (d) -- (p); \draw[ar] (d) -- (dp);
  \node[draw=amber!35, fill=amberbr!7, rounded corners=4pt, minimum width=10.2cm,
        minimum height=0.55cm, text=ink2] at (4.6,-1.4)
    {~Tiga masukan, dua keluaran, satu simpul antara yang dipakai bersama --- tak mungkin ditulis sebagai Sequential.};
\end{tikzpicture}

}
\end{center}
\vskip 3pt{\color{ink3}\fontsize{7.4}{9.4}\selectfont Tiga masukan (judul, badan teks, tag), dua keluaran (skor prioritas dan departemen).\par}
\icode{python}{listing 7.9}{listings/ch07-03.py}
\begin{iband}[signal]
Buku menyebutnya \hilite{seperti bermain LEGO}: cara yang sederhana tetapi sangat luwes untuk mendefinisikan graf lapis yang sebarang.
\end{iband}
\end{frame}

\framekicker{Bagian 7.2.2 $\cdot$ listing 7.10-7.11}
\begin{frame}[shrink=25]{Melatihnya: senarai berurutan, atau kamus bernama}
\begin{minipage}[t]{0.485\textwidth}
\icode{python}{listing 7.10 --- senarai}{listings/ch07-04.py}

\end{minipage}%
\hfill
\begin{minipage}[t]{0.485\textwidth}
\icode{python}{listing 7.11 --- kamus}{listings/ch07-05.py}

\end{minipage}%
\begin{iband}[signal]
Versi senarai \textbf{harus mengikuti urutan} yang Anda berikan ke konstruktor \inlinecode{Model()}. Kalau masukan atau keluarannya banyak, \hilite{pakai kamus bernama} -- urutannya jadi tidak penting lagi, dan kodenya jauh lebih tahan salah.
\end{iband}
\end{frame}

\framekicker{Bagian 7.2.2 $\cdot$ listing 7.12-7.13}
\begin{frame}[shrink=25]{Kekuatan sesungguhnya: akses ke keterhubungan lapis}
{\fontsize{9.4}{12.6}\selectfont\color{fg2}Model Functional adalah \textbf{struktur data graf yang eksplisit}. Itu memungkinkan dua hal yang tidak bisa dilakukan model subclass: \textbf{penggambaran topologi} dan \textbf{ekstraksi fitur}.\par}\vskip 4pt
\begin{minipage}[t]{0.485\textwidth}
\icode{python}{menggambar topologi}{listings/ch07-06.py}

\end{minipage}%
\hfill
\begin{minipage}[t]{0.485\textwidth}
\icode{python}{listing 7.12 --- memeriksa simpul}{listings/ch07-07.py}

\end{minipage}%
\icode{python}{listing 7.13 --- ekstraksi fitur}{listings/ch07-08.py}
\begin{iband}[signal]
\inlinecode{None} pada shape tensor menyatakan \textbf{ukuran batch} -- model ini menerima batch berukuran berapa pun.
\end{iband}
\end{frame}

\framekicker{Bagian 7.2.3 $\cdot$ listing 7.14}
\begin{frame}[shrink=25]{Model subclassing: kendali penuh}
\icode{python}{listing 7.14}{listings/ch07-09.py}
\begin{itable}{>{\raggedright\arraybackslash}p{0.3196\textwidth}>{\raggedright\arraybackslash}p{0.3102\textwidth}>{\raggedright\arraybackslash}p{0.3102\textwidth}}
\thead{} & \thead{\inlinecode{Layer}} & \thead{\inlinecode{Model}} \\ \midrule
Perannya & Bata bangunan untuk membuat model. & Objek tingkat atas yang benar-benar dilatih dan diekspor. \\
\inlinecode{fit()}, \inlinecode{evaluate()}, \inlinecode{predict()} & Tidak punya. & Punya. \\
Disimpan ke berkas & Tidak. & Bisa. \\
\end{itable}
\begin{iband}[signal]
Selebihnya kedua kelas itu \textbf{praktis identik}. Subclassing membuka model yang \hilite{tidak bisa dinyatakan sebagai graf berarah tanpa siklus} -- misalnya \inlinecode{call()} yang memakai lapis di dalam gelung \inlinecode{for}, atau bahkan memanggilnya secara rekursif.
\end{iband}
\end{frame}

\framekicker{Bagian 7.2.3}
\begin{frame}[shrink=25]{Harga kebebasan itu}
\begin{minipage}[t]{0.485\textwidth}
{\fontsize{9.4}{12.6}\selectfont\color{fg2}\textbf{Model Functional} = struktur data eksplisit -- graf lapis yang bisa Anda lihat, periksa, dan ubah.\par}\vskip 4pt
{\fontsize{9.4}{12.6}\selectfont\color{fg2}\textbf{Model subclass} = \hilite{sepotong bytecode} -- kelas Python dengan metode \inlinecode{call()} berisi kode mentah. Itulah sumber keluwesannya, dan sekaligus sumber batasannya.\par}\vskip 4pt

\end{minipage}%
\hfill
\begin{minipage}[t]{0.485\textwidth}
\begin{minipage}[t]{1.0\textwidth}
\begin{icardbad}
\cardh{Yang hilang saat subclassing}
\inlinecode{summary()} \textbf{tidak menampilkan keterhubungan lapis} $\cdot$ \inlinecode{plot\_model()} \textbf{tidak bisa dipakai} $\cdot$ \textbf{ekstraksi fitur tidak mungkin} -- sebab memang tidak ada grafnya.
\end{icardbad}
\end{minipage}%

\end{minipage}%
\begin{iband}[rose]
Begitu model subclass diinstansiasi, \textbf{lintasan majunya menjadi kotak hitam sepenuhnya}. Anda mengembangkan objek Python baru, bukan sekadar menjepretkan bata LEGO -- \hilite{permukaan galatnya jauh lebih luas, dan pekerjaan awakutunya lebih banyak}.
\end{iband}
\end{frame}

\framekicker{Bagian 7.2.4-7.2.5 $\cdot$ listing 7.15-7.16}
\begin{frame}[shrink=25]{Ketiganya bisa dicampur -- dan mana yang dipilih}
\begin{minipage}[t]{0.485\textwidth}
\icode{python}{listing 7.15 --- subclass di dalam Functional}{listings/ch07-10.py}

\end{minipage}%
\hfill
\begin{minipage}[t]{0.485\textwidth}
\icode{python}{listing 7.16 --- Functional di dalam subclass}{listings/ch07-11.py}

\end{minipage}%
\begin{iband}[signal]
\textbf{Saran buku:} kalau model Anda bisa dinyatakan sebagai graf berarah tanpa siklus, \hilite{pakai Functional API, bukan subclassing}. Seluruh contoh di sisa buku memakai Functional API -- tetapi dengan \textbf{lapis-lapis subclass} di dalamnya. Kombinasi itu memberi keluwesan pengembangan sekaligus keuntungan Functional API.
\end{iband}
\end{frame}

\coursesection{02}{Lingkar pelatihan bawaan}{Metrik kustom, callback, dan TensorBoard.}

\framekicker{Bagian 7.3.1 $\cdot$ listing 7.18}
\begin{frame}[shrink=25]{Menulis metrik sendiri}
{\fontsize{9.4}{12.6}\selectfont\color{fg2}Metrik Keras adalah subkelas \inlinecode{keras.metrics.Metric}. Seperti lapis, ia punya \textbf{keadaan internal} yang disimpan di variabel Keras. Bedanya, variabel itu \hilite{tidak diperbarui lewat backpropagation}, jadi Anda menulis sendiri logika pembaruannya.\par}\vskip 4pt
\icode{python}{listing 7.18 --- RMSE sebagai metrik kustom}{listings/ch07-12.py}
\begin{iband}[signal]
Tiga metode itulah kontraknya: \textbf{\inlinecode{update\_state()}} memperbarui, \textbf{\inlinecode{result()}} melaporkan nilai sekarang, \textbf{\inlinecode{reset\_state()}} mengosongkan tanpa perlu membuat objek baru.
\end{iband}
\end{frame}

\framekicker{Bagian 7.3.2}
\begin{frame}[shrink=25]{Callback: dari pesawat kertas jadi dron}
\begin{iquote}
\quotetext{Meluncurkan pelatihan pada data besar selama puluhan epoch dengan \inlinecode{model.fit()} itu sedikit seperti melempar pesawat kertas: setelah dorongan awal, Anda tidak punya kendali atas lintasannya maupun tempat mendaratnya. API callback Keras mengubah panggilan \inlinecode{model.fit()} Anda dari pesawat kertas menjadi \textbf{dron cerdas dan otonom} yang bisa memeriksa dirinya sendiri dan bertindak.}\quotecite{Chollet \& Watson, bab 7.3.2}
\end{iquote}
\begin{minipage}[t]{0.241\textwidth}
\begin{icard}
\cardh{Model checkpointing}
Menyimpan keadaan model di berbagai titik selama pelatihan.
\end{icard}
\end{minipage}%
\hfill
\begin{minipage}[t]{0.241\textwidth}
\begin{icard}
\cardh{Early stopping}
Menghentikan pelatihan saat rugi validasi berhenti membaik -- dan menyimpan model terbaik yang sudah diperoleh.
\end{icard}
\end{minipage}%
\hfill
\begin{minipage}[t]{0.241\textwidth}
\begin{icard}
\cardh{Menyetel parameter dinamis}
Misalnya learning rate optimizer, di tengah pelatihan.
\end{icard}
\end{minipage}%
\hfill
\begin{minipage}[t]{0.241\textwidth}
\begin{icard}
\cardh{Mencatat \& menggambar}
Bilah kemajuan \inlinecode{fit()} yang sudah Anda kenal itu \hilite{sebenarnya sebuah callback}
\end{icard}
\end{minipage}%
\icode{python}{callback bawaan (tidak lengkap)}{listings/ch07-13.py}
\end{frame}

\framekicker{Bagian 7.3.2 $\cdot$ listing 7.19}
\begin{frame}[shrink=25]{EarlyStopping + ModelCheckpoint: pasangan baku}
\icode{python}{listing 7.19}{listings/ch07-14.py}
\begin{iband}[signal]
Ini menggantikan pola boros di bab 4-5: latih sampai overfit untuk \textbf{mencari} epoch terbaik, lalu latih ulang dari nol sebanyak itu. Dengan pasangan callback ini, \hilite{sekali jalan sudah cukup}
\end{iband}
\icode{python}{menyimpan dan memuat manual}{listings/ch07-15.py}
\end{frame}

\framekicker{Bagian 7.3.3 $\cdot$ listing 7.20}
\begin{frame}[shrink=25]{Callback kustom -- enam kait yang tersedia}
\begin{minipage}[t]{0.485\textwidth}
\icode{python}{kait yang bisa Anda isi}{listings/ch07-16.py}
{\fontsize{9.4}{12.6}\selectfont\color{fg2}Semuanya dipanggil dengan argumen \inlinecode{logs}, sebuah kamus berisi informasi tentang batch, epoch, atau jalannya pelatihan -- metrik latih dan validasi.\par}\vskip 4pt

\end{minipage}%
\hfill
\begin{minipage}[t]{0.485\textwidth}
\icode{python}{listing 7.20 --- riwayat rugi per batch}{listings/ch07-17.py}

\end{minipage}%
\end{frame}

\framekicker{Bagian 7.3.4}
\begin{frame}[shrink=25]{TensorBoard: memutar lingkar kemajuan lebih cepat}
{\fontsize{9.4}{12.6}\selectfont\color{fg2}Kemajuan adalah proses berulang: \textbf{gagasan $\rightarrow$ percobaan $\rightarrow$ hasil $\rightarrow$ gagasan berikutnya} (gambar 7.6). Makin banyak putaran yang bisa Anda jalankan, makin tajam gagasannya. Keras memperpendek jarak gagasan-ke-percobaan; GPU cepat memperpendek percobaan-ke-hasil; \hilite{TensorBoard menangani hasil-ke-gagasan berikutnya}\par}\vskip 4pt
\begin{minipage}[t]{0.241\textwidth}
\begin{icard}
\cardh{Pantau metrik}
Secara visual, selama pelatihan berjalan.
\end{icard}
\end{minipage}%
\hfill
\begin{minipage}[t]{0.241\textwidth}
\begin{icard}
\cardh{Gambar arsitektur}
Visualisasi model.
\end{icard}
\end{minipage}%
\hfill
\begin{minipage}[t]{0.241\textwidth}
\begin{icard}
\cardh{Histogram}
Aktivasi dan gradien.
\end{icard}
\end{minipage}%
\hfill
\begin{minipage}[t]{0.241\textwidth}
\begin{icard}
\cardh{Embedding 3D}
Jelajahi ruang embedding.
\end{icard}
\end{minipage}%
\icode{python}{memakainya}{listings/ch07-18.py}
\icode{bash}{menjalankan servernya}{listings/ch07-19.sh}
\end{frame}

\coursesection{03}{Menulis lingkar pelatihan sendiri}{Saat fit() tidak cukup -- dan cara tetap memakainya.}

\framekicker{Bagian 7.4}
\begin{frame}[shrink=25]{Kapan fit() memang tidak cukup}
{\fontsize{9.4}{12.6}\selectfont\color{fg2}\inlinecode{fit()} bawaan \textbf{hanya berfokus pada supervised learning}: ada target yang diketahui, dan rugi dihitung sebagai fungsi target itu dan prediksi model. Tidak semua bentuk machine learning masuk kategori itu.\par}\vskip 4pt
\begin{minipage}[t]{0.3253\textwidth}
\begin{icard}
\cardh{Generative learning}
Tidak ada target eksplisit. Diperkenalkan di \textbf{bab 16}.
\end{icard}
\end{minipage}%
\hfill
\begin{minipage}[t]{0.3253\textwidth}
\begin{icard}
\cardh{Self-supervised}
Targetnya diambil \textbf{dari masukannya sendiri}.
\end{icard}
\end{minipage}%
\hfill
\begin{minipage}[t]{0.3253\textwidth}
\begin{icard}
\cardh{Reinforcement learning}
Pembelajaran didorong \textbf{imbalan sesekali} -- seperti melatih anjing.
\end{icard}
\end{minipage}%
\begin{iband}[amber]
Dua kehalusan yang harus diperhatikan saat menulis lingkar sendiri:
\end{iband}
\begin{isteps}
\item \textbf{\inlinecode{training=True} wajib diteruskan.} Lapis seperti \inlinecode{Dropout} berperilaku berbeda saat latih dan saat inferensi. \inlinecode{dropout(inputs, training=True)} menjatuhkan sebagian aktivasi; \inlinecode{training=False} tidak melakukan apa pun. Jadi: \inlinecode{predictions = model(inputs, training=True)}.
\item \textbf{Pakai \inlinecode{model.trainable\_weights}, bukan \inlinecode{model.weights}.} Model punya dua jenis bobot: \textbf{terlatih} (diperbarui backpropagation) dan \textbf{tak terlatih} (diperbarui lapis pemiliknya saat lintasan maju). Di antara lapis bawaan Keras, satu-satunya yang punya bobot tak terlatih adalah \textbf{\inlinecode{BatchNormalization}} -- ia melacak rerata dan simpangan baku data yang lewat (bab 9).
\end{isteps}
\end{frame}

\framekicker{Bagian 7.4.2}
\begin{frame}[shrink=25]{Satu langkah pelatihan, tiga backend}
\icode{python}{TensorFlow}{listings/ch07-20.py}
\icode{python}{PyTorch}{listings/ch07-21.py}
\begin{iband}[rose]
Pada PyTorch, \inlinecode{model.zero\_grad()} \textbf{kritis}: panggilan \inlinecode{backward()} bersifat menambahkan. Kalau gradien tidak dikosongkan tiap langkah, nilainya menumpuk dan \hilite{pelatihan tidak akan berjalan}
\end{iband}
\end{frame}

\framekicker{Bagian 7.4.2 $\cdot$ JAX}
\begin{frame}[shrink=25]{JAX: paling rumit, karena tanpa keadaan sama sekali}
{\fontsize{9.4}{12.6}\selectfont\color{fg2}Karena fungsi gradiennya diperoleh lewat metapemrograman, Anda harus lebih dulu menuliskan fungsi yang \textbf{mengembalikan} rugi. Fungsi itu harus tanpa keadaan: ia menerima semua variabel yang dipakainya sebagai argumen, dan \textbf{mengembalikan nilai setiap variabel yang diperbaruinya} -- termasuk bobot tak terlatih tadi.\par}\vskip 4pt
\icode{python}{stateless\_call() dan value\_and\_grad}{listings/ch07-22.py}
\begin{iband}[signal]
Pola yang sama berlaku untuk metrik: di JAX, metode yang mengubah keadaan seperti \inlinecode{update\_state()} tidak bisa dipanggil di dalam fungsi tanpa keadaan. Padanannya: \hilite{\inlinecode{stateless\_update\_state()}, \inlinecode{stateless\_result()}, \inlinecode{stateless\_reset\_state()}}
\end{iband}
\end{frame}

\framekicker{Bagian 7.4.3}
\begin{frame}[shrink=25]{Memakai metrik di tingkat rendah}
\begin{minipage}[t]{0.485\textwidth}
\icode{python}{metrik biasa}{listings/ch07-23.py}
\ioutput{listings/ch07-24.txt}

\end{minipage}%
\hfill
\begin{minipage}[t]{0.485\textwidth}
\icode{python}{melacak rerata skalar}{listings/ch07-25.py}
\ioutput{listings/ch07-26.txt}

\end{minipage}%
\begin{iband}[amber]
Jangan lupa \inlinecode{metric.reset\_state()} saat mau mengosongkan hasilnya -- \textbf{di awal tiap epoch pelatihan, dan di awal evaluasi}. Lupa melakukannya membuat angka metrik \hilite{tercampur antar-epoch}
\end{iband}
\end{frame}

\framekicker{Bagian 7.4.4 $\cdot$ listing 7.21}
\begin{frame}[shrink=25]{Jalan tengah: train\_step() sendiri, fit() tetap dipakai}
{\fontsize{9.4}{12.6}\selectfont\color{fg2}Menulis seluruh lingkar dari nol memberi keluwesan paling besar, tetapi Anda kehilangan \textbf{callback, pengoptimalan unjuk kerja, dan dukungan pelatihan tersebar}. Jalan tengahnya: ganti \inlinecode{train\_step()} saja, biarkan kerangka mengerjakan sisanya.\par}\vskip 4pt
\icode{python}{listing 7.21 --- versi TensorFlow}{listings/ch07-27.py}
\begin{ibullets}
\item Pola ini \textbf{tidak menghalangi} Anda memakai Functional API -- berlaku untuk Sequential, Functional, maupun subclass.
\item Anda \textbf{tidak perlu} dekorator \inlinecode{@tf.function} atau \inlinecode{@jax.jit}; \hilite{kerangkanya yang mengerjakan itu untuk Anda}.
\end{ibullets}
\end{frame}

\framekicker{Ringkasan}
\begin{frame}[shrink=25]{Yang wajib terbawa dari bab 7}
\begin{isteps}
\item \textbf{Progressive disclosure of complexity} -- satu spektrum alur kerja di atas API bersama \inlinecode{Layer} dan \inlinecode{Model}, bukan beberapa kerangka terpisah.
\item \textbf{Sequential} untuk tumpukan sederhana; \textbf{Functional API} untuk graf lapis -- dan itu yang dipakai sisa buku ini; \textbf{subclassing} untuk yang tidak bisa dinyatakan sebagai graf.
\item Functional memberi \textbf{akses ke keterhubungan lapis}: \inlinecode{plot\_model()} dan ekstraksi fitur. Subclassing \hilite{membuang keduanya}.
\item Campuran terbaik: \textbf{model Functional yang berisi lapis-lapis subclass}.
\item \textbf{Metrik kustom} = \inlinecode{update\_state} / \inlinecode{result} / \inlinecode{reset\_state}. \textbf{Callback} = enam kait \inlinecode{on\_*}.
\item \textbf{\inlinecode{EarlyStopping} + \inlinecode{ModelCheckpoint}} menggantikan pola latih-ulang yang boros itu.
\item Menulis \inlinecode{train\_step()} sendiri: ingat \textbf{\inlinecode{training=True}} dan \textbf{\inlinecode{trainable\_weights}}; di PyTorch jangan lupa \textbf{\inlinecode{zero\_grad()}}; di JAX semuanya lewat padanan \textbf{\inlinecode{stateless\_*}}.
\end{isteps}
\vskip 4pt
\reslink{NOTEBOOK}{03\_train\_step\_kustom\_per\_backend.ipynb}{../../course-slides/notebooks/ch07/03\_train\_step\_kustom\_per\_backend.ipynb}
\reslink{BAB BERIKUT}{Bab 8 --- Klasifikasi citra}{../ch08/index.html}
\vskip 2pt
\end{frame}

\end{document}
